\documentclass[10pt,a4paper]{amsart}
\usepackage[margin=19mm]{geometry}
\usepackage{amsmath,amssymb,mathtools}
\usepackage[scr=rsfs,cal=euler]{mathalfa}
\usepackage{xfrac}
\usepackage[unicode,hidelinks,bookmarks=false]{hyperref}
\newcommand{\ie}{i.e.\ }
\newcommand{\cf}{cf.\ }
\newcommand{\resp}{respectively\ }
\newcommand{\Subsection}{\subsection}
\providecommand{\nobreakdash}{\nobreak\hbox{-}}
\setlength{\emergencystretch}{2em}
\multlinegap0pt
\usepackage{xcolor}
\usepackage{amssymb}
\usepackage{mathtools}
\usepackage{enumerate}
\newtheorem{thm}{Theorem}
\title{Uniqueness for embeddings of nuclear $C^*$-algebras into type II$_{1}$ factors}
\author{Shanshan Hua, Stuart White}
\date{Source: arXiv:2601.08779v2 (CC BY 4.0)}
\begin{document}
\maketitle

\noindent\emph{Source and licence.} Uniqueness for embeddings of nuclear $C^*$-algebras into type II$_{1}$ factors. Version arXiv:2601.08779v2 (CC BY 4.0), distributed by arXiv under the Creative Commons Attribution 4.0 International licence, http://creativecommons.org/licenses/by/4.0/, as recorded in the arXiv licence metadata for this exact version and shown on the abstract page https://arxiv.org/abs/2601.08779v2. Full licence notice: \texttt{licenses/arxiv-2601.08779-CC-BY-4.0.txt}.\par\smallskip

\noindent\textit{Editorial scope.} Counted unit: the article's thm Kirchberg (source0.tex lines 1373--1375). The article's complete original proof of it is delivered with it (source0.tex lines 1384--1388, one \texttt{proof} environment of 5 lines, which the article places away from the statement). No mathematics is written, repaired or re-derived: the statement and the proof are the article's own lines, in the article's own notation, and no retained line is substituted or re-flowed.  Retained uncounted as the source-local context the proof needs: lines 1368--1370.  Nothing was omitted from any retained span, so the package carries no omission record.  No retained line contains a citation, so the wrapper carries no bibliography and no bibliography entry.  No retained line contains a figure environment, an image-inclusion command or drawing code, so no image is delivered and the package declares no asset dependency.  Editorial formatting only, in addition: an AMS \texttt{amsart} preamble that restates the article's own declarations and macros used by the retained lines, namely the environments \texttt{thm} on separate counters, together with the standard packages those lines need.  The retained environments print in this document as Theorem 1 in order of appearance, whereas the article prints the same items with the numbering of its own full document.  The complete native source of the article as distributed in the arXiv e-print for this exact version is bundled unchanged as \texttt{sources/192625/source0.tex}.
\begin{thm}\label{Appendix:Thm}
Let $A$ be a separable, unital and exact $C^*$-algebra, and let $\mathcal M$ be a type III factor with separable predual. Then any two unital and nuclear $^*$-homomorphisms $A\to\mathcal M$ are approximately unitarily equivalent if and only if they have the same kernel.
\end{thm}

\begin{thm}[Kirchberg]\label{Kirchberg}
Let $A$ be a separable, unital and exact $C^*$-algebra, and let $B$ be a unital, simple purely infinite $C^*$-algebra.  Then two unital, injective and nuclear $^*$-homomorphisms $\phi,\psi:A\to B$ are asymptotically unitarily equivalent (i.e. there is a continuous path of unitaries $(u_t)_{t\geq 0}$ in $B$ with $u_t\phi(a)u_t^*\to \psi(a)$ as $t\to\infty$ for all $a\in A$) if and only if $KK_{\mathrm{nuc}}(\phi)=KK_{\mathrm{nuc}}(\psi)$.
\end{thm}

\begin{proof}[Proof of Theorem \ref{Appendix:Thm}]
Certainly if $\phi$ and $\psi$ are approximately unitarily equivalent, they have the same kernel. Conversely if they have the same kernel $I$, we can replace $A$ by $A/I$ and it suffices to show that two unital, nuclear and injective maps $\phi,\psi:A\to\mathcal M$ are approximately or even asymptotically unitarily equivalent, which will follow from Theorem \ref{Kirchberg} by showing $KK_{\mathrm{nuc}}(\phi)=0$.

Fix a family $(v_n)_n$ of isometries in $\mathcal M$ with pairwise orthogonal ranges such that $\sum_{n=1}^\infty v_nv_n^*=1$ (with strong convergence). Then given a unital and nuclear $^*$-homomorphism $\phi:A\to\mathcal M$, one can define the infinite repeat $\phi^{(\infty)}:A\to \mathcal M$ by $\phi(x)=\sum_{n=1}^\infty v_n\phi(x)v_n^*$. By construction we will have $[\phi^{(\infty)}]+[\phi]=[\phi^{(\infty)}]$ in $KK_{\mathrm{nuc}}(A,\mathcal M)$, and hence $[\phi]=0$ in $KK_{\mathrm{nuc}}(A,\mathcal M)$.
\end{proof}

\end{document}
